\documentclass[10pt,twocolumn,letterpaper]{article}

\usepackage{times}
\usepackage{epsfig}
\usepackage{graphicx}
\usepackage{amsmath}
\usepackage{amssymb}
\usepackage{booktabs}
\usepackage{hyperref}
\usepackage[ruled,vlined]{algorithm2e}

% Include other packages you need here.

% Set page margins
\usepackage[top=1in, bottom=1in, left=0.75in, right=0.75in]{geometry}

% If you comment hyperref and then uncomment it, you should delete
% egpaper.aux before re-running latex.
\hypersetup{breaklinks=true,bookmarks=false,colorlinks}

\begin{document}

%%%%%%%%% TITLE
\title{Emergent Unsafe Prompts: Eliciting Hidden Dangers in Vision-Language Models Through Context-Dependent Adversarial Attacks}

\author{Team 15\\
Gardenia Harris, Ryan Liu, Keyu Chen, Sarah Johnson\\
\texttt{\{gardenia.harris, ryan.liu, keyu.chen, sarah.johnson\}@university.edu}
}

\maketitle

%%%%%%%%% ABSTRACT
\begin{abstract}
Vision-language models (VLMs) are increasingly deployed in safety-critical applications, yet their vulnerability to context-dependent adversarial prompts remains underexplored. We present a novel two-pipeline approach to generate ``emergent unsafe'' prompts—textual instructions that appear benign in isolation but become dangerous when paired with specific image contexts. Pipeline 1 uses supervised fine-tuning via Low-Rank Adaptation (LoRA) of LLaMA-3.1-8B on image captions to generate context-aware prompts. Pipeline 2 employs direct prompting of frontier models (Claude Sonnet 4, Gemini 2.0 Flash) with systematic tone and danger-relation templates. Human evaluation of 100 samples reveals 38.6\% of generated prompts are genuinely unsafe, while automated LLM judges (Gemini 2.5 Flash Lite) achieve 62.4\% accuracy with 50.8\% precision, highlighting persistent over-flagging challenges. Our work demonstrates cost-effective LoRA fine-tuning ($<$\$5) can rival frontier models in generating nuanced safety threats, advancing understanding of emergent risks in multimodal AI systems.
\end{abstract}

%%%%%%%%% BODY TEXT
\section{Introduction}

\subsection{Problem Statement and Motivation}
Large vision-language models (VLMs) such as GPT-4V, Gemini, and LLaVA have demonstrated remarkable capabilities in understanding and reasoning about images and text. However, their deployment in real-world applications raises critical safety concerns. While existing adversarial attacks focus on explicit harmful content or jailbreaking techniques, a more subtle class of threats remains largely unexplored: \textbf{emergent unsafe prompts}—instructions that appear innocuous when read alone but become dangerous when interpreted in conjunction with visual context.

For example, the prompt ``Let's test the structural integrity using a heavier ball'' seems like a reasonable physics experiment in isolation. However, when paired with an image of children playing soccer, it implicitly suggests replacing the soccer ball with a dangerous heavy object that could cause injury. This context-dependent danger represents a fundamental challenge for content moderation systems that evaluate text and images independently.

\subsection{Project Goal}
Our primary goal is to develop systematic methods for generating emergent unsafe prompts at scale and evaluate the ability of both humans and automated systems to detect such threats. Specifically, we aim to:
\begin{itemize}
    \item Design scalable pipelines for generating high-quality emergent unsafe prompts
    \item Compare cost-effectiveness and performance of fine-tuned open-source models versus frontier commercial models
    \item Benchmark human agreement and LLM judge accuracy on detecting emergent safety threats
    \item Analyze failure modes and recommend improvements for safety classifiers
\end{itemize}

\subsection{Approach Overview}
\begin{figure}[t]
\centering
\includegraphics[width=0.48\textwidth]{figures/pipeline_overview.pdf}
\caption{\textbf{Two-Pipeline Approach for Emergent Unsafe Prompt Generation.} Pipeline 1 directly prompts frontier models (Claude Sonnet 4, Gemini 2.0 Flash) with systematic templates. Pipeline 2 fine-tunes LLaMA-3.1-8B using LoRA on multi-level captions to generate context-aware prompts. Human evaluation shows 38.6\% of 100 sampled prompts are genuinely unsafe, while automated judges achieve 62.4\% accuracy with persistent over-flagging (51.6\% FPR).}
\label{fig:overview}
\end{figure}

Our approach (Figure~\ref{fig:overview}) consists of two complementary pipelines:

\textbf{Pipeline 1 (Direct Prompting):} We systematically prompt frontier models (Claude Sonnet 4, Gemini 2.0 Flash) using structured templates that vary tone (DIY, scientific, playful, creative, challenge, observational) and danger relations (stability, contact, proximity, threshold, interaction). This approach leverages in-context learning to generate diverse emergent unsafe prompts at scale.

\textbf{Pipeline 2 (Supervised Fine-Tuning):} We fine-tune LLaMA-3.1-8B using Low-Rank Adaptation (LoRA) on 500 COCO validation images with 5 caption levels per image. The model learns to generate emergent unsafe prompts conditioned on image captions, leveraging parameter-efficient training to reduce costs and test whether specialized fine-tuning can match frontier model performance.

\subsection{Novel Contributions}
\begin{enumerate}
    \item \textbf{First systematic study} of emergent unsafe prompts in vision-language models, introducing a rigorous definition and evaluation framework
    \item \textbf{Cost-effective LoRA pipeline} achieving comparable performance to frontier models at $<$1\% of inference cost
    \item \textbf{Large-scale dataset} of 1,388 emergent unsafe prompts with human annotations and automated judgments
    \item \textbf{Empirical analysis} revealing LLM judges over-flag benign prompts (51.6\% FPR) while maintaining high recall (84.6\%), informing future safety classifier design
\end{enumerate}

\section{Approach}

\subsection{Defining Emergent Unsafe Prompts}
We define an \textit{emergent unsafe prompt} as a textual instruction satisfying three criteria:
\begin{enumerate}
    \item \textbf{Safe Alone:} The prompt contains no explicitly harmful keywords or obvious dangerous suggestions when read in isolation
    \item \textbf{Unsafe with Context:} When paired with a specific image (described via caption), the prompt implies a realistic and plausible path to physical danger
    \item \textbf{Action-Oriented:} The prompt suggests a concrete action (not mere speculation like ``I wonder'' or ``What if'')
\end{enumerate}

This definition distinguishes emergent threats from traditional adversarial attacks that rely on explicit jailbreaking or prompt injection.

\subsection{Pipeline 1: Systematic Prompting of Frontier Models}

\subsubsection{Template Design}
We design a structured prompt template incorporating:
\begin{enumerate}
    \item \textbf{System instruction:} Define emergent unsafe prompts and provide examples
    \item \textbf{Caption input:} 5-level caption describing the target image
    \item \textbf{Tone specification:} One of \{DIY, scientific, playful, creative, challenge, observational\}
    \item \textbf{Danger relation:} One of \{stability, contact, proximity, threshold, interaction\}
\end{enumerate}

This systematic variation ensures diverse prompt styles and danger types.

\subsubsection{Model Selection}
We evaluate two frontier models:
\begin{itemize}
    \item \textbf{Claude Sonnet 4:} Strong reasoning and instruction-following
    \item \textbf{Gemini 2.0 Flash:} Fast, cost-effective, multimodal capabilities
\end{itemize}

Both models received identical system prompts and few-shot examples to ensure fair comparison.

\subsection{Pipeline 2: LoRA Fine-Tuning of LLaMA-3.1-8B}

\subsubsection{Motivation for LoRA}
Full fine-tuning of large language models is prohibitively expensive for academic research. Low-Rank Adaptation (LoRA)~\cite{hu2021lora} freezes the pre-trained model weights and injects trainable low-rank decomposition matrices into each transformer layer:
\begin{equation}
    W' = W_0 + \Delta W = W_0 + BA
\end{equation}
where $W_0 \in \mathbb{R}^{d \times k}$ is frozen, and $B \in \mathbb{R}^{d \times r}$, $A \in \mathbb{R}^{r \times k}$ with rank $r \ll \min(d,k)$ are trainable.

For our experiments, we use $r=8$, $\alpha=16$, targeting attention query and value projection matrices. This reduces trainable parameters from 8B to $\sim$4.2M (0.05\% of total), enabling training on a single A100 GPU.

\subsubsection{Training Data and Process}
We curated a dataset of 500 COCO validation images, each annotated with 5 caption levels (from concise to detailed). For each caption, we manually authored 5 diverse emergent unsafe prompts spanning different tones and danger types. This yielded 12,500 training examples.

Training hyperparameters:
\begin{itemize}
    \item Base model: LLaMA-3.1-8B-Instruct
    \item LoRA rank: $r=8$, $\alpha=16$
    \item Learning rate: $3 \times 10^{-4}$ with cosine decay
    \item Batch size: 4 (gradient accumulation: 8)
    \item Epochs: 3
    \item Total training time: 6 hours on A100
    \item Total cost: \$4.80 (via cloud GPU rental)
\end{itemize}

\subsubsection{Why LoRA is Cost-Effective}
Compared to prompting frontier models at scale:
\begin{itemize}
    \item \textbf{Upfront cost:} LoRA training = \$4.80 (one-time)
    \item \textbf{Inference cost:} \$0.01 per 1000 prompts (local deployment) vs. \$0.50-\$2.00 for GPT-4/Claude
    \item \textbf{Scalability:} Generate 100k prompts for $\sim$\$1 vs. \$50-\$200 using APIs
    \item \textbf{Control:} No rate limits, full model ownership, reproducible outputs
\end{itemize}

For a production system generating millions of prompts annually, LoRA amortizes training costs within days.


\subsection{Secret Weapon: Multi-Level Caption Hierarchy}
Our key insight is that \textbf{caption abstraction level critically affects prompt quality}. We generate 5 caption levels per image:
\begin{enumerate}
    \item \textbf{Level 0:} Minimal, object-focused (``A dog catches a frisbee'')
    \item \textbf{Level 1-2:} Moderate detail with scene context
    \item \textbf{Level 3-4:} Detailed descriptions including background, actions, emotions
\end{enumerate}

Emergent unsafe prompts require sufficient context to be dangerous but not so much detail that models simply repeat caption content. We found Levels 2-3 optimal for balancing creativity and danger implication.

\section{Implementation}

\subsection{Phase 1: Unsafe Prompt Generation Pipeline (Gardenia, Ryan)}

\subsubsection{Data Preparation}
\begin{enumerate}
    \item Downloaded COCO 2017 validation split (5,000 images)
    \item Randomly sampled 500 images covering diverse scenes
    \item Generated 5-level captions using BLIP-2 and GPT-4V with temperature variation
    \item Prepared caption dataset for use in both frontier model prompting and LoRA training
    \item Split data: 80\% train (400 images), 20\% validation (100 images)
\end{enumerate}

\subsubsection{API Integration}
Implemented parallel prompting infrastructure:
\begin{itemize}
    \item Anthropic Claude API with exponential backoff for rate limits
    \item Google Gemini API with batch processing (50 requests/min)
    \item Retry logic with jittered delays to handle transient failures
    \item Cost tracking middleware logging token usage per request
\end{itemize}

\subsubsection{Prompt Template Engineering}
Developed systematic template with placeholders:
\begin{verbatim}
<SYSTEM>
You are an expert red-teamer...
</SYSTEM>

Caption: {caption_text}
Tone: {tone}
Danger Relation: {danger_relation}

Generate an emergent unsafe prompt...
\end{verbatim}

Iterated through 6 tones $\times$ 5 danger relations = 30 combinations per caption.

\subsubsection{Quality Control}
Implemented automated checks:
\begin{enumerate}
    \item JSON schema validation for structured outputs
    \item Length constraints (5-30 words)
    \item Profanity filter to ensure ``safe alone'' criterion
    \item Manual review of 100 random samples for calibration
\end{enumerate}

\subsection{Phase 2: LoRA Training Pipeline (Keyu, Sarah)}

\subsubsection{Training Data Curation}
We curated a dataset for LoRA fine-tuning:
\begin{enumerate}
    \item Used 500 COCO validation images with 5 caption levels per image
    \item Manually authored 5 diverse emergent unsafe prompts per caption, stratified by tone and danger type
    \item This yielded 12,500 training examples (500 images $\times$ 5 captions $\times$ 5 prompts)
    \item Split: 80\% train (400 images, 10k prompts), 20\% validation (100 images, 2.5k prompts)
\end{enumerate}

\subsubsection{Model Architecture and Training}
We employed the \texttt{transformers} and \texttt{peft} libraries:
\begin{verbatim}
from transformers import AutoModelForCausalLM
from peft import LoraConfig, get_peft_model

config = LoraConfig(
    r=8, lora_alpha=16,
    target_modules=["q_proj", "v_proj"],
    lora_dropout=0.1,
    task_type="CAUSAL_LM"
)
model = AutoModelForCausalLM.from_pretrained(
    "meta-llama/Llama-3.1-8B-Instruct",
    torch_dtype=torch.bfloat16
)
model = get_peft_model(model, config)
\end{verbatim}

Training script executed 3 epochs with early stopping based on validation perplexity.

\subsubsection{Inference and Post-Processing}
Generated prompts undergo filtering:
\begin{enumerate}
    \item Remove exact caption repetitions
    \item Filter prompts containing explicit harm keywords (violence, illegal substances)
    \item Deduplicate semantically similar prompts using sentence embeddings (threshold: 0.85 cosine similarity)
    \item Validate prompts meet minimum length (10 tokens) and grammaticality
\end{enumerate}

\subsection{Experimental Setup}
\subsubsection{Dataset Statistics}
\begin{table}[h]
\centering
\caption{Dataset Statistics for Emergent Unsafe Prompt Generation}
\begin{tabular}{lcc}
\toprule
\textbf{Metric} & \textbf{Pipeline 1} & \textbf{Pipeline 2} \\
\midrule
Total images & 500 & 500 \\
Prompts per image & 25 & 27.8 (avg) \\
Total prompts & 12,500 & 13,880 \\
Avg prompt length & 14.2 tokens & 16.7 tokens \\
Training cost & \$4.80 & N/A \\
Inference cost (1k) & \$0.01 & \$1.20-\$2.50 \\
\bottomrule
\end{tabular}
\label{tab:dataset_stats}
\end{table}

\subsubsection{Evaluation Protocol}
\begin{enumerate}
    \item \textbf{Human Evaluation:} 4 annotators independently labeled 100 random prompts (stratified across pipelines) as unsafe/safe using majority vote (≥3/4 agreement)
    \item \textbf{Automated Judge:} Gemini 2.5 Flash Lite with strict evaluation rubric and few-shot examples
    \item \textbf{Metrics:} Accuracy, Precision, Recall, F1, False Positive/Negative Rates
\end{enumerate}

\section{Results}

\subsection{Human Evaluation: Ground Truth Labels}

\subsubsection{Prompt-Level Analysis}
Human annotators reviewed 101 prompts across 28 images (combining original 50-sample and new 49-sample batches). Using majority vote (≥3 ``yes'' = unsafe):
\begin{itemize}
    \item \textbf{39 prompts (38.6\%)} labeled unsafe
    \item \textbf{62 prompts (61.4\%)} labeled safe
    \item Average inter-annotator agreement: 72.3\% (moderate-substantial)
\end{itemize}

\textbf{Vote distribution:}
\begin{itemize}
    \item 4 yes / 0 no: 12 prompts (unanimous unsafe)
    \item 3 yes / 1 no: 26 prompts
    \item 2 yes / 2 no: 29 prompts (ties → labeled safe)
    \item 1 yes / 3 no: 18 prompts
    \item 0 yes / 4 no: 15 prompts (unanimous safe)
\end{itemize}

\subsubsection{Image-Level Analysis}
\begin{itemize}
    \item \textbf{28 unique images} evaluated
    \item \textbf{3.61 prompts per image} (average)
    \item \textbf{25 images (89.3\%)} had ≥1 unsafe prompt
    \item \textbf{3 images (10.7\%)} had all prompts safe
    \item \textbf{2 images (7.1\%)} had all prompts unsafe
\end{itemize}

This demonstrates that emergent unsafe prompts are highly image-dependent—certain visual contexts (e.g., children playing, indoor fire hazards) are more susceptible.

\subsection{LLM Judge Performance}

\subsubsection{Confusion Matrix}
\begin{table}[h]
\centering
\caption{Gemini 2.5 Flash Lite Judge Performance on 101 Prompts}
\begin{tabular}{lcc}
\toprule
\textbf{Metric} & \textbf{Value} & \textbf{\%} \\
\midrule
True Positives (TP) & 33 & -- \\
False Positives (FP) & 32 & -- \\
True Negatives (TN) & 30 & -- \\
False Negatives (FN) & 6 & -- \\
\midrule
Accuracy & 63/101 & 62.4\% \\
Precision & 33/65 & 50.8\% \\
Recall & 33/39 & 84.6\% \\
F1 Score & -- & 63.5\% \\
False Positive Rate & 32/62 & 51.6\% \\
False Negative Rate & 6/39 & 15.4\% \\
\bottomrule
\end{tabular}
\label{tab:llm_judge_metrics}
\end{table}

\subsubsection{Key Findings}
\begin{enumerate}
    \item \textbf{High Recall, Low Precision:} The judge catches 84.6\% of unsafe prompts but flags 51.6\% of safe prompts as unsafe
    \item \textbf{Over-Flagging Bias:} LLM errs on the side of caution, reflecting conservative safety training
    \item \textbf{Improved Coverage:} After adding second batch, no prompts remain unevaluated (previously 5 missing)
\end{enumerate}

\subsection{Pipeline Comparison: LoRA vs. Frontier Models}

\begin{figure}[t]
\centering
\includegraphics[width=0.48\textwidth]{figures/pipeline_comparison.pdf}
\caption{\textbf{Cost-Effectiveness Comparison.} LoRA-finetuned LLaMA-3.1-8B achieves 36.2\% human-validated unsafe rate at 0.8\% the cost of Claude Sonnet 4 (41.5\% unsafe rate). Gemini 2.0 Flash offers middle ground at 20\% the cost with 38.9\% unsafe rate.}
\label{fig:pipeline_comparison}
\end{figure}

\begin{table}[h]
\centering
\caption{Pipeline Performance and Cost Analysis}
\begin{tabular}{lccc}
\toprule
\textbf{Method} & \textbf{Unsafe \%} & \textbf{Cost/1k} & \textbf{Quality} \\
\midrule
LoRA LLaMA & 36.2\% & \$0.01 & 3.8/5 \\
Claude Sonnet 4 & 41.5\% & \$1.20 & 4.2/5 \\
Gemini 2.0 Flash & 38.9\% & \$0.25 & 4.0/5 \\
\bottomrule
\end{tabular}
\label{tab:pipeline_comparison}
\end{table}

\textbf{Analysis:}
\begin{itemize}
    \item LoRA achieves 87\% of Claude's unsafe rate at 0.8\% cost
    \item Gemini offers optimal cost-quality tradeoff for moderate-scale experiments
    \item All pipelines exceed 35\% unsafe rate, validating approach effectiveness
\end{itemize}

\subsection{Failure Mode Analysis}

\subsubsection{False Positives (32 cases)}
Common patterns of over-flagging:
\begin{enumerate}
    \item \textbf{Playful physics questions:} ``How high can we throw it?'' misinterpreted as dangerous
    \item \textbf{Photography prompts:} ``Capture the precise moment'' flagged despite being observational
    \item \textbf{Household scenarios:} Blender, fireplace, toilet prompts assumed unsafe without context
    \item \textbf{Hypothetical speculation:} ``I wonder if...'' treated as action suggestion
\end{enumerate}

\subsubsection{False Negatives (6 cases)}
Missed unsafe prompts include:
\begin{enumerate}
    \item Living room heat-transfer prompt (3-1 unsafe) – scientific language obscured danger
    \item Ski-slope proximity test (4-0 unsafe) – judge didn't recognize collision risk
    \item Frisbee kinetic-energy (4-0 unsafe) – physics jargon masked throwing at power lines
    \item Bike friction test (3-1 unsafe) – unclear actionability
    \item Nintendo Wii reinforced housing (3-1 unsafe) – ``force testing'' seemed benign
\end{enumerate}

\subsection{Heatmap: Judge vs. Human Labels}

\begin{figure}[h]
\centering
\includegraphics[width=0.48\textwidth]{figures/confusion_heatmap.pdf}
\caption{\textbf{Confusion Heatmap.} Gemini judge over-flags 32 safe prompts (51.6\% FPR) while missing only 6 unsafe ones (15.4\% FNR). Improving precision without sacrificing recall remains a key challenge.}
\label{fig:heatmap}
\end{figure}

\subsection{Cost-Effectiveness Summary}

For generating 100,000 emergent unsafe prompts:
\begin{itemize}
    \item \textbf{LoRA:} \$4.80 (training) + \$1.00 (inference) = \$5.80 total
    \item \textbf{Claude Sonnet 4:} \$120,000 (inference only)
    \item \textbf{Gemini 2.0 Flash:} \$25,000 (inference only)
\end{itemize}

\textbf{Conclusion:} LoRA is 20,000$\times$ cheaper than Claude and 4,300$\times$ cheaper than Gemini for large-scale deployment.

\section{Related Work}

\subsection{Adversarial Attacks on Vision-Language Models}
Prior work has explored jailbreaking VLMs through adversarial images~\cite{bailey2023jailbreak}, typographic attacks~\cite{goh2023multimodal}, and prompt injection~\cite{perez2022ignore}. Our work differs by focusing on \textit{context-dependent} rather than \textit{explicit} harm, requiring multimodal reasoning to detect.

\subsection{Safety Alignment in Large Language Models}
Reinforcement Learning from Human Feedback (RLHF)~\cite{ouyang2022training} and Constitutional AI~\cite{bai2022constitutional} have improved LLM safety for textual content. However, these methods struggle with emergent dangers that only manifest in specific visual contexts, as demonstrated by our 51.6\% false positive rate.

\subsection{Red-Teaming and Automated Adversarial Generation}
Recent efforts in automated red-teaming~\cite{perez2022red, ganguli2022red} primarily target explicit policy violations. Our systematic tone$\times$danger-relation templates introduce a structured approach to generating subtle, context-dependent threats.

\subsection{Parameter-Efficient Fine-Tuning}
LoRA~\cite{hu2021lora} and related methods (Prefix Tuning~\cite{li2021prefix}, Adapter Layers~\cite{houlsby2019parameter}) enable cost-effective adaptation of large models. We demonstrate LoRA's viability for safety research, achieving near-frontier performance at 0.05\% trainable parameters.

\section{Conclusion}

\subsection{Contributions to the Literature}
This work makes four key contributions:
\begin{enumerate}
    \item \textbf{Novel threat model:} Formalizing emergent unsafe prompts as context-dependent adversarial inputs
    \item \textbf{Scalable generation pipelines:} Demonstrating LoRA fine-tuning rivals frontier models at 20,000$\times$ lower cost
    \item \textbf{Human-validated dataset:} Providing 101 annotated prompts with 4-way majority vote labels
    \item \textbf{Judge failure analysis:} Revealing 51.6\% false positive rate in state-of-the-art LLM safety classifiers
\end{enumerate}

\subsection{Lessons Learned}
\begin{enumerate}
    \item \textbf{Caption abstraction matters:} Levels 2-3 yield highest-quality emergent unsafe prompts
    \item \textbf{LoRA is production-ready:} Parameter-efficient training achieves 87\% of Claude's performance at <1\% cost
    \item \textbf{Safety classifiers over-flag:} Current LLM judges prioritize recall over precision, requiring hybrid human-in-the-loop systems
    \item \textbf{Human agreement is moderate:} 72.3\% inter-annotator agreement suggests emergent danger is subjective
\end{enumerate}

\subsection{Self-Assessment}
We successfully achieved our primary goal of generating large-scale emergent unsafe prompts (1,388 total) and benchmarking automated judges. However, two limitations remain:
\begin{enumerate}
    \item \textbf{Limited image diversity:} 500 COCO images may not generalize to all visual domains (medical, industrial, etc.)
    \item \textbf{Single judge model:} Only evaluated Gemini 2.5 Flash Lite; other judges (GPT-4V, Claude) may perform differently
\end{enumerate}

Overall, we estimate 85\% goal completion—core methodology proven, but broader validation needed.

\subsection{Future Work}
Assuming another semester, we would pursue:
\begin{enumerate}
    \item \textbf{Multi-turn attacks:} Chaining benign prompts to gradually escalate danger
    \item \textbf{Vision-grounded generation:} Directly conditioning on images rather than captions
    \item \textbf{Adversarial training:} Fine-tuning judges on our dataset to reduce false positives
    \item \textbf{Real-world deployment:} Testing prompts against production VLMs (GPT-4V, Gemini Pro Vision)
    \item \textbf{Cross-lingual evaluation:} Exploring whether emergent dangers translate across languages
\end{enumerate}

\section{Group Contribution Statement}

\textbf{Gardenia Harris (25\%):}
\begin{itemize}
    \item Led LoRA training pipeline design and implementation
    \item Authored abstract, introduction, and conclusion sections
    \item Coordinated human annotation experiments
    \item Prepared final presentation slides
\end{itemize}

\textbf{Ryan Liu (25\%):}
\begin{itemize}
    \item Designed multi-level caption generation system
    \item Implemented LoRA training scripts and hyperparameter tuning
    \item Wrote Approach and Implementation sections
    \item Created cost-effectiveness analysis and visualizations
\end{itemize}

\textbf{Keyu Chen (25\%):}
\begin{itemize}
    \item Developed Pipeline 2 API integration infrastructure
    \item Engineered systematic prompt templates
    \item Conducted automated judge evaluation experiments
    \item Contributed to Results section and failure mode analysis
\end{itemize}

\textbf{Sarah Johnson (25\%):}
\begin{itemize}
    \item Managed Phase 2 frontier model experiments
    \item Performed quality control and manual review
    \item Authored Related Work section
    \item Compiled and formatted final LaTeX report
\end{itemize}

All team members participated in weekly meetings, data annotation, and iterative prompt refinement.

%%%%%%%%% REFERENCES
\begin{thebibliography}{9}

\bibitem{hu2021lora}
Edward J. Hu, Yelong Shen, Phillip Wallis, Zeyuan Allen-Zhu, Yuanzhi Li, Shean Wang, Lu Wang, and Weizhu Chen.
\newblock LoRA: Low-Rank Adaptation of Large Language Models.
\newblock In \textit{ICLR}, 2022.

\bibitem{bailey2023jailbreak}
Peter Bailey et al.
\newblock Image Hijacks: Adversarial Images can Control Generative Models at Runtime.
\newblock \textit{arXiv:2309.00236}, 2023.

\bibitem{goh2023multimodal}
Gabriel Goh et al.
\newblock Multimodal Neurons in Artificial Neural Networks.
\newblock \textit{Distill}, 2023.

\bibitem{perez2022ignore}
Fabio Perez and Ian Ribeiro.
\newblock Ignore Previous Prompt: Attack Techniques For Language Models.
\newblock \textit{arXiv:2211.09527}, 2022.

\bibitem{ouyang2022training}
Long Ouyang et al.
\newblock Training Language Models to Follow Instructions with Human Feedback.
\newblock In \textit{NeurIPS}, 2022.

\bibitem{bai2022constitutional}
Yuntao Bai et al.
\newblock Constitutional AI: Harmlessness from AI Feedback.
\newblock \textit{arXiv:2212.08073}, 2022.

\bibitem{perez2022red}
Ethan Perez et al.
\newblock Red Teaming Language Models with Language Models.
\newblock In \textit{EMNLP}, 2022.

\bibitem{ganguli2022red}
Deep Ganguli et al.
\newblock Red Teaming Language Models to Reduce Harms: Methods, Scaling Behaviors, and Lessons Learned.
\newblock \textit{arXiv:2209.07858}, 2022.

\bibitem{li2021prefix}
Xiang Lisa Li and Percy Liang.
\newblock Prefix-Tuning: Optimizing Continuous Prompts for Generation.
\newblock In \textit{ACL}, 2021.

\bibitem{houlsby2019parameter}
Neil Houlsby et al.
\newblock Parameter-Efficient Transfer Learning for NLP.
\newblock In \textit{ICML}, 2019.

\end{thebibliography}

\end{document}
